\documentclass[12pt]{article}
\usepackage{amsfonts,amsthm,amsmath,amssymb,mathtools}
\usepackage{mathrsfs}
\usepackage{enumitem}
\usepackage{tikz-cd}
\usepackage{geometry}
\usepackage{mlmodern}

\geometry{letterpaper, portrait, margin=1in}

\setcounter{section}{0}

\renewcommand{\thesection}{\S\arabic{section}}
\renewcommand{\thesubsection}{\arabic{section}}
\newcommand{\dd}{\mathop{}\!\mathrm{d}}

\newtheorem{thm}{Theorem}
\newenvironment{theorem}[1]{
	\renewcommand{\thethm}{#1}
	\begin{thm}
		}{\end{thm}}

\newtheorem{lma}{Lemma}
\newenvironment{lemma}[1]{
	\renewcommand{\thelma}{#1}
	\begin{lma}
		}{\end{lma}}

\theoremstyle{definition}
\newtheorem{ex}{}
\newenvironment{exercise}[1]{
	\renewcommand{\theex}{#1}
	\begin{ex}
		}{\end{ex}}

\title{Homework 8}
\author{Connor Mooneyhan}
\date{October 16, 2025}

\begin{document}

\begin{center}
	{\bf\Large MTH 346/646 homework cover sheet}
\end{center}

\vspace{0.5in}

\noindent{\large Name}: Connor Mooneyhan \\
{\large Homework number}: 8 \\
{\large Date}: October 23, 2025 \\

{\Large What topics did this assignment cover? How comfortable do you feel with
each of those topics?}

It covered finite generation of curves, or the lack thereof.
I feel quite comfortable with the theory.


\vspace{1.75in}

{\Large In the course of doing this assignment, where did you get stuck?
	How did you get unstuck?}

  I wouldn't say I got particularly stuck at any point, though I did have to try a few different things to figure out how to force 2(b).


\vspace{1.75in}

{\Large Is there anything you've done here that you're not confident about, that you want me to take a really close look at, or that you want me to suggest an alternate approach to?}

I'm really not sure how to find the second generator in 1, or even how to prove that it can be generated by only 2 elements.

\pagebreak

\begin{ex}
	Let $E : y^{2} = x^{3} + x^{2} - 2x + 1$. The goal of this problem is to use Magma
	to find two points $P$ and $Q$ so that $E(\mathbb{Q})$ is generated by $P$ and $Q$. (You may use Magma for computations, but don't use the commands {\tt Generators}
	or {\tt MordellWeilGroup}).

	The information below gives you explicit forms of the four lemmas.

	\begin{enumerate}[label=(\alph*)]
		\item You may assume that $[E(\mathbb{Q}) : 2E(\mathbb{Q})] = 2$ and coset representatives are
		      $Q_{1} = (0 : 1 : 0)$ and $Q_{2} = (-2,1)$.
		\item You may assume that if $P_{0} \in E(\mathbb{Q})$, then for all $P \in E(\mathbb{Q})$ we have
		      $h(P_{0} + P) \leq 2h(P) + 2h(P_{0}) + 3.8821$.
		      Also, for all $P \in E(\mathbb{Q})$, $h(2P) \geq 4h(P) - 3.8821$.
		\item The Magma command {\tt NaiveHeight} of a point $P$ returns $h(P)$.
		      Also, the command \verb+Points(E : Bound := B)+ returns all points $P$
		      with $h(P) \leq \log(B)$.
	\end{enumerate}
\end{ex}
\begin{proof}
	% Since $E(\mathbb{Q})$ is abelian, $2E(\mathbb{Q})$ is a normal subgroup, so since it has index 2, \begin{equation*}
	%   E(\mathbb{Q})/2E(\mathbb{Q}) \cong \mathbb{Z}/2 \mathbb{Z}.
	% \end{equation*}
	% The plan is to find a generating 
	We first want to find, in the language of the Descent Theorem, our $\kappa'$ and $\kappa$.
	The statement of (b) above tells us immediately that $\kappa = 3.8821$.
	Also by (b), $\kappa_i$ for a given $P_i \in E(\mathbb{Q})$ is $2h(P_i) + 3.8821$, so we want to take the max of the $\kappa_i$'s corresponding to the $Q_i$'s.
	With this notation, Magma tells us that $\kappa_1 = 3.8821$ and $\kappa_2 \approx 5.2684$.
	Thus we can define $\kappa' = 5.2684$.
	The number we most want out of this is $\kappa + \kappa' = 9.1505$.

	We use this to conclude that if we have a $P \in E(\mathbb{Q})$ such that $h(P) \geq 9.1505$, then we can find an $R \in E(\mathbb{Q})$ through the process outlined in the proof of the Descent Theorem such that $h(R) \leq 9.1505$ and
	\begin{equation*}
		P = a_1Q_1 + a_2Q_2 + 2^mR = a_2Q_2 + 2^mR
	\end{equation*}
	for some $a_1, a_2, m \in \mathbb{Z}$, where the $a_1Q_1$ term disappears because $Q_1$ is the identity, so $a_1Q_1$ is also the identity.
	Since $e^{9.1505} \approx 9419.1488$, we can find all such possibilities for $R$ by running \begin{verbatim}
  Points(E : Bound := 9419);.
  \end{verbatim}
	This gives us many points, namely \begin{verbatim}
  (0 : 1 : 0), (0 : 1 : 1), (0 : -1 : 1), (1 : 1 : 1), (1 : -1 : 1),
  (-2 : 1 : 1), (-2 : -1 : 1), (2 : 3 : 1), (2 : -3 : 1), (-3/4 : 13/8 : 1),
  (-3/4 : -13/8 : 1), (4/9 : 17/27 : 1), (4/9 : -17/27 : 1), (12 : 43 : 1),
  (12 : -43 : 1), (-55/49 : 603/343 : 1), (-55/49 : -603/343 : 1),
  (70/121 : 811/1331 : 1), (70/121 : -811/1331 : 1), (154/25 : 2017/125 : 1),
  (154/25 : -2017/125 : 1), (2193/2704 : 106079/140608 : 1),
  (2193/2704 : -106079/140608 : 1), (-4472/2601 : 201709/132651 : 1),
  (-4472/2601 : -201709/132651 : 1), (5304/1849 : 413867/79507 : 1),
  (5304/1849 : -413867/79507 : 1). \end{verbatim}
	That's 27 points.
	I can tell you for certain that, assuming that I did my calculations right with $\kappa$ and $\kappa'$, then $Q_2$ together with all the points listed above does indeed generate $E(\mathbb{Q})$, and thus it is finitely generated.
	Showing, however, that it is indeed generated by two elements, seems a bit more daunting.
	I suppose I would have to find a point that generates the subgroup $2E(\mathbb{Q})$, which I'm not even sure is cyclic in the first place.
	It stands to reason that if any point in $E(\mathbb{Q})$ generates all 27 of these points, then it is probably one of these points, since heights only tend to increase.
	I tried narrowing down which points I had to test by finding the order of each one with the \verb+Order()+ function, but that only eliminated three possibilities.

	I decided to brute force it at this point and see if any of these points generated all the others using the following code, expecting a final count of 27 for a generator: \begin{verbatim}
    E := EllipticCurve([0, 1, 0, -2, 1]);
    Q1 := E![0, 1, 0];
    Q2 := E![-2, 1];
    points := Points(E : Bound := 9419);

    for point in points do;
      count := 0;
      for i in [-50..50] do;
        if i * point in points then;
          count := count + 1;
        end if;
      end for;
      printf "%o: %o\n", point, count;
    end for; \end{verbatim}
	The best this got me was a few points that generate 9 of the 27, but nothing better than that.
	So that being said, I'm not sure which of those points I should choose to accompany $Q_2$ as the second generator of $E(\mathbb{Q})$.
	(One of the points that generated 9 of them was $(1 : 1 : 1)$ in case that counts for something.)
\end{proof}

\begin{ex}
	$*$ Let $T : x^{2} + y^{2} = 1$ and
	$T(\mathbb{Q}) = \{ (x,y) : x, y \in \mathbb{Q} \text{ and } x^{2} + y^{2} = 1
		\}$. (We don't include points at infinity in $T$.)

	The set $T(\mathbb{Q})$ has the structure of a group, with the group
	operation defined by
	\[
		(x_{1},y_{1}) * (x_{2},y_{2}) = (x_{1} x_{2} - y_{1} y_{2}, x_{1} y_{2} + x_{2} y_{1}).
	\]
	(This is equivalent to adding angles, and the formula above is probably
	reminiscent of the addition formulas for cosine and sine.) Note: You may assume that this operation gives $T(\mathbb{Q})$ the structure of a group. You don't need to prove it.

	\begin{enumerate}[label=(\alph*)]
		\item Let $S = \{ p_{1}, \ldots, p_{k} \}$ be a finite set of prime numbers.
		      Show that if $Q_{1}$ and $Q_{2}$ are points in $T(\mathbb{Q})$ where the denominator
		      of the $x$ and $y$ coordinates have only prime factors from $S$,
		      then $Q_{1} * Q_{2}$ has the same property.
		\item Show that if $S$ is any finite set of prime numbers there is
		      some point $(x,y) \in T(\mathbb{Q})$ where the denominators of $x$ and $y$ contain
		      prime factors not in $S$. (The parametrization of rational points on $x^{2} + y^{2} = 1$ might be helpful here!)
		\item Conclude that $T(\mathbb{Q})$ is not finitely generated. (The main way
		      the descent theorem breaks in this case is because $[T(\mathbb{Q}) : 2T(\mathbb{Q})]$ is
		      not finite.)
	\end{enumerate}
\end{ex}
\begin{proof}\
	\begin{enumerate}[label=(\alph*)]
		\item Write \begin{equation*}
			      Q_1 = \left(\frac{a}{p_1^{a_1}\cdots p_k^{a_k}}, \frac{b}{p_1^{b_1}\cdots p_k^{b_k}}\right), Q_2 = \left(\frac{c}{p_1^{c_1}\cdots p_k^{c_k}}, \frac{d}{p_1^{d_1}\cdots p_k^{d_k}}\right)
		      \end{equation*}
		      for some integers $a, b, c, d, a_1, \dots, a_k, b_1, \dots, b_k, c_1, \dots, c_k, d_1, \dots, d_k$.
		      Then \begin{align*}
			       & Q_1 * Q_2                                                                                                                                                                                                                                                                                                                                          \\
			       & = \left( \frac{ac}{p_1^{a_1c_1} \cdots p_k^{a_kc_k}} - \frac{bd}{p_1^{b_1d_1} \cdots p_k^{b_kd_k}}, \frac{ad}{p_1^{a_1d_1} \cdots p_k^{a_kd_k}} + \frac{bc}{p_1^{b_1c_1} \cdots p_k^{b_kc_k}} \right)                                                                                                                                              \\
			       & = \left( \frac{ac\left(p_1^{b_1d_1} \cdots p_k^{b_kd_k}\right) - bd\left(p_1^{a_1c_1} \cdots p_k^{a_kc_k}\right)}{p_1^{a_1c_1 + b_1d_1} \cdots p_k^{a_kc_k + b_kd_k}}, \frac{ad \left( p_1^{b_1c_1} \cdots p_k^{b_kc_k} \right) + bc \left( p_1^{a_1d_1} \cdots p_k^{a_kd_k} \right)}{p_1^{a_1d_1 + b_1c_1} \cdots p_k^{a_kd_k + b_kc_k}} \right),
		      \end{align*}
		      which has the desired property.
        \item Let $S$ and $Q_1$ have notation as in (a).
          Let $m = \alpha/\beta$, where $\alpha, \beta \in \mathbb{Z}$ and $m$ is not the slope of the line tangent to $x^2 + y^2 = 1$ at $Q_1$.
      We can then find the unique point $P = (X, Y) \in T(\mathbb{Q})$ such that the line between $Q_1$ and $P$ has slope $m$.

      First, we find the $y$-intercept of this line, which will be \begin{equation*}
        B = \frac{b}{p_1^{b_1}\cdots p_k^{b_k}} - \frac{a\alpha}{\beta p_1^{a_1} \cdots p_k^{a_k}} = \frac{b\beta p_1^{a_1} \cdots p_k^{a_k} - a\alpha p_1^{b_1} \cdots p_k^{b_k}}{\beta p_1^{a_1 + b_1} \cdots p_k^{a_k + b_k}}.
      \end{equation*}
      Thus \begin{align*}
        Y & = \frac{\alpha X}{\beta} + \frac{b\beta p_1^{a_1} \cdots p_k^{a_k} - a\alpha p_1^{b_1} \cdots p_k^{b_k}}{\beta p_1^{a_1 + b_1} \cdots p_k^{a_k + b_k}} \\
        & = \frac{\alpha X p_1^{a_1 + b_1} \cdots p_k^{a_k + b_k} + b\beta p_1^{a_1} \cdots p_k^{a_k} - a\alpha p_1^{b_1} \cdots p_k^{b_k}}{\beta p_1^{a_1 + b_1} \cdots p_k^{a_k + b_k}}.
      \end{align*}
      Certainly, we can choose $\beta$ such that it contains no factors in $S$, and then $Y$ will be rational with denominator not consisting of prime factors only in $S$, as desired.
		\item Consider the subgroup $G$ (which is a subgroup by (a)) generated by some finite set of points $P_1, \dots, P_m \in T(\mathbb{Q})$.
		      The set of primes represented in the denominators of $P_1, \dots, P_m$ is a finite set (call it $S$), and every linear combination of $P_1, \dots, P_m$ (i.e. every element of $G$) has a denominator which is a product of primes in $S$.
		      By (b), then, there exist points of $T(\mathbb{Q})$ which are not in $G$.
		      Since $G$ was an arbitrary finitely generated subgroup, this shows that every finitely generated subgroup of $T(\mathbb{Q})$ is a proper subgroup, and thus $T(\mathbb{Q})$ is not finitely generated.
	\end{enumerate}
\end{proof}

\end{document}
